\documentclass[10pt,a4paper]{amsart}
\usepackage{amsmath,amssymb,amsthm,graphicx}
% Retained preamble for the Stacks Project chapter morphisms.tex (Morphisms of Schemes).
% The chapter file pinned-original-source.tex carries only \input{preamble} on its first line, so no
% preamble text lives inside the chapter itself.  The Stacks Project uses one single shared preamble for
% every chapter; the snapshot stages it verbatim in sources/stacks-chow/preamble.tex.  The block below is
% copied mechanically, line for line, from that file's lines 14-18 (the Xy-pic setup) and lines 154-217
% (the native theorem declarations together with the mathematical macros).  Nothing was added, removed or
% redefined.  The wrapper class and the amsmath/amssymb/amsthm/graphicx base packages are supplied by the
% assembler, so the class/\externaldocument part of the shared preamble is not repeated here.
% For commutative diagrams we use Xy-pic
\usepackage[all]{xy}

% We use 2cell for 2-commutative diagrams.
\xyoption{2cell}
% Theorem environments.
%
\theoremstyle{plain}
\newtheorem{theorem}[subsection]{Theorem}
\newtheorem{proposition}[subsection]{Proposition}
\newtheorem{lemma}[subsection]{Lemma}

\theoremstyle{definition}
\newtheorem{definition}[subsection]{Definition}
\newtheorem{example}[subsection]{Example}
\newtheorem{exercise}[subsection]{Exercise}
\newtheorem{situation}[subsection]{Situation}

\theoremstyle{remark}
\newtheorem{remark}[subsection]{Remark}
\newtheorem{remarks}[subsection]{Remarks}

\numberwithin{equation}{subsection}

% Macros
%
\def\lim{\mathop{\mathrm{lim}}\nolimits}
\def\colim{\mathop{\mathrm{colim}}\nolimits}
\def\Spec{\mathop{\mathrm{Spec}}}
\def\Hom{\mathop{\mathrm{Hom}}\nolimits}
\def\Ext{\mathop{\mathrm{Ext}}\nolimits}
\def\SheafHom{\mathop{\mathcal{H}\!\mathit{om}}\nolimits}
\def\SheafExt{\mathop{\mathcal{E}\!\mathit{xt}}\nolimits}
\def\Sch{\mathit{Sch}}
\def\Mor{\mathop{\mathrm{Mor}}\nolimits}
\def\Ob{\mathop{\mathrm{Ob}}\nolimits}
\def\Sh{\mathop{\mathit{Sh}}\nolimits}
\def\NL{\mathop{N\!L}\nolimits}
\def\CH{\mathop{\mathrm{CH}}\nolimits}
\def\proetale{{pro\text{-}\acute{e}tale}}
\def\etale{{\acute{e}tale}}
\def\QCoh{\mathit{QCoh}}
\def\Ker{\mathop{\mathrm{Ker}}}
\def\Im{\mathop{\mathrm{Im}}}
\def\Coker{\mathop{\mathrm{Coker}}}
\def\Coim{\mathop{\mathrm{Coim}}}

% Boxtimes
%
\DeclareMathSymbol{\boxtimes}{\mathbin}{AMSa}{"02}

%
% Macros for moduli stacks/spaces
%
\def\QCohstack{\mathcal{QC}\!\mathit{oh}}
\def\Cohstack{\mathcal{C}\!\mathit{oh}}
\def\Spacesstack{\mathcal{S}\!\mathit{paces}}
\def\Quotfunctor{\mathrm{Quot}}
\def\Hilbfunctor{\mathrm{Hilb}}
\def\Curvesstack{\mathcal{C}\!\mathit{urves}}
\def\Polarizedstack{\mathcal{P}\!\mathit{olarized}}
\def\Complexesstack{\mathcal{C}\!\mathit{omplexes}}
% \Pic is the operator that assigns to X its picard group, usage \Pic(X)
% \Picardstack_{X/B} denotes the Picard stack of X over B
% \Picardfunctor_{X/B} denotes the Picard functor of X over B
\def\Pic{\mathop{\mathrm{Pic}}\nolimits}
\def\Picardstack{\mathcal{P}\!\mathit{ic}}
\def\Picardfunctor{\mathrm{Pic}}
\def\Deformationcategory{\mathcal{D}\!\mathit{ef}}
\usepackage{cleveref}
\crefname{originalsection}{original section}{original sections}
\Crefname{originalsection}{Original section}{Original sections}
\makeatletter
\newlabel{schemes-section-quasi-coherent}{{01LA}{source}{}{sourcepointer-1}{}}
\newlabel{schemes-section-quasi-coherent@cref}{{[originalsection][0][]01LA}{[1][0][]source}{}{}{}}
\makeatother
\title{A submodule of a quasi-coherent module contains a largest quasi-coherent submodule, characterised by its universal property}
\author{The Stacks Project Authors}
\date{stacks-project:morphisms.tex:local-snapshot-20261012; GNU FDL 1.2 or later}
\begin{document}
\maketitle
\noindent\textit{Editorial note.} All delivered mathematical text is the author's own, copied from the pinned Stacks Project chapter file morphisms.tex; no statement or proof step was rewritten, reordered or removed.\par\medskip
\noindent\textit{Editorial note.} The canonical Stacks tag of the counted statement is 01QZ (morphisms-lemma-largest-quasi-coherent-subsheaf). Following the multi-claim convention for this source, source\_url is the source-level URL https://stacks.math.columbia.edu/ and the tag is carried in claim\_key and here rather than in a per-claim tag URL.\par\medskip
\noindent\textit{Editorial note.} The preamble is the chapter's shared Stacks Project preamble, retained as sources/stacks-morphisms/retained-preamble.tex: a line-for-line excerpt of sources/stacks-chow/preamble.tex lines 14-18 and 154-217. No macro was invented or redefined.\par\medskip
\noindent\textit{Editorial note.} The statement is self-contained (F, G, G' and the universal property are all named in it), so no context span is attached; the two same-chapter lemmas that surround it in the source are separate claims and are not counted here.\par\medskip
\noindent\textit{Allowed standard prerequisites (editorial).}\begin{itemize}
\item Images and direct sums of maps of quasi-coherent O\_X-modules are quasi-coherent. Schemes, Section labelled schemes-section-quasi-coherent (canonical Stacks tag 01LA); invoked at proof lines 465-468. Supplies the two closure properties that make the constructed image sheaf G' quasi-coherent, and again makes the image of H -> G -> F quasi-coherent.
\item A quasi-coherent O\_X-submodule of a quasi-coherent O\_X-module is a kernel of a map of quasi-coherent modules, hence itself quasi-coherent, and the set of quasi-coherent submodules contained in a fixed submodule is a set. Background from the same chapter and from Schemes, Section schemes-section-quasi-coherent (canonical Stacks tag 01LA). Makes the indexing set A of all quasi-coherent submodules contained in G available so that the direct sum over A can be formed.
\end{itemize}\medskip
\ifdefined\hypertarget\hypertarget{sourcepointer-1}{}\fi
\noindent\textit{Source reference.} Original section 01LA: sources/stacks-chow/tags.txt, source lines 2079--2079. Cross-chapter pointer invoked at proof lines 465-468 for the two closure properties (direct sums and images of maps of quasi-coherent sheaves) on which the construction rests; the construction and the universal property are carried out in full inside the counted span.\par\medskip
\begin{lemma}
\label{lemma-largest-quasi-coherent-subsheaf}
Let $X$ be a scheme. Let $\mathcal{F}$ be a quasi-coherent
$\mathcal{O}_X$-module. Let $\mathcal{G} \subset \mathcal{F}$
be a $\mathcal{O}_X$-submodule. There exists a unique quasi-coherent
$\mathcal{O}_X$-submodule $\mathcal{G}' \subset \mathcal{G}$
with the following property: For every quasi-coherent $\mathcal{O}_X$-module
$\mathcal{H}$ the map
$$
\Hom_{\mathcal{O}_X}(\mathcal{H}, \mathcal{G}')
\longrightarrow
\Hom_{\mathcal{O}_X}(\mathcal{H}, \mathcal{G})
$$
is bijective. In particular $\mathcal{G}'$ is the largest quasi-coherent
$\mathcal{O}_X$-submodule of $\mathcal{F}$ contained in $\mathcal{G}$.
\end{lemma}

\begin{proof}
Let $\mathcal{G}_a$, $a \in A$ be the set of quasi-coherent
$\mathcal{O}_X$-submodules contained in $\mathcal{G}$.
Then the image $\mathcal{G}'$ of
$$
\bigoplus\nolimits_{a \in A} \mathcal{G}_a \longrightarrow \mathcal{F}
$$
is quasi-coherent as the image of a map of quasi-coherent sheaves
on $X$ is quasi-coherent and since a direct sum of quasi-coherent sheaves
is quasi-coherent, see
Schemes, Section \ref{schemes-section-quasi-coherent}.
The module $\mathcal{G}'$ is contained in $\mathcal{G}$. Hence this is the
largest quasi-coherent $\mathcal{O}_X$-module contained in $\mathcal{G}$.

\medskip\noindent
To prove the formula, let $\mathcal{H}$ be a quasi-coherent
$\mathcal{O}_X$-module and let $\alpha : \mathcal{H} \to \mathcal{G}$
be an $\mathcal{O}_X$-module map. The image of the composition
$\mathcal{H} \to \mathcal{G} \to \mathcal{F}$ is quasi-coherent
as the image of a map of quasi-coherent sheaves. Hence it is contained
in $\mathcal{G}'$. Hence $\alpha$ factors through $\mathcal{G}'$
as desired.
\end{proof}

\end{document}
